\chapter{What D'Artagnan Went to Paris For}

\lettrine{T}{he} lieutenant dismounted before a shop in the Rue des Lombards, at the sign of the \buildingname{Pilon d'Or.} A man of good appearance, wearing a white apron, and stroking his gray moustache with a large hand, uttered a cry of joy on perceiving the pied horse. <Monsieur le chevalier,> said he, <ah, is that you?>

<Bon jour, Planchet,> replied D'Artagnan, stooping to enter the shop.

<Quick, somebody,> cried Planchet, <to look after Monsieur d'Artagnan's horse,—somebody to get ready his room,—somebody to prepare his supper.>

<Thanks, Planchet. Good-day, my children!> said D'Artagnan to the eager boys.

<Allow me to send off this coffee, this treacle, and these raisins,> said Planchet; <they are for the store-room of monsieur le surintendant.>

<Send them off, send them off!>

<That is only the affair of a moment, then we shall sup.>

<Arrange it that we may sup alone; I want to speak to you.>

Planchet looked at his old master in a significant manner.

<Oh, don't be uneasy, it is nothing unpleasant,> said D'Artagnan.

<So much the better—so much the better!> And Planchet breathed freely again, whilst D'Artagnan seated himself quietly down in the shop, upon a bale of corks, and made a survey of the premises. The shop was well stocked; there was a mingled perfume of ginger, cinnamon, and ground pepper, which made D'Artagnan sneeze. The shop-boy, proud of being in company with so renowned a warrior, of a lieutenant of musketeers, who approached the person of the king, began to work with an enthusiasm which was something like delirium, and to serve the customers with a disdainful haste that was noticed by several.

Planchet put away his money, and made up his accounts, amidst civilities addressed to his former master. Planchet had with his equals the short speech and haughty familiarity of the rich shopkeeper who serves everybody and waits for nobody. D'Artagnan observed this habit with a pleasure which we shall analyse presently. He saw night come on by degrees, and at length Planchet conducted him to a chamber on the first story, where, amidst bales and chests, a table very nicely set out awaited the two guests.

D'Artagnan took advantage of a moment's pause to examine the countenance of Planchet, whom he had not seen for a year. The shrewd Planchet had acquired a slight protuberance in front, but his countenance was not puffed. His keen eye still played with facility in its deep-sunk orbit; and fat, which levels all the characteristic saliences of the human face, had not yet touched either his high cheek-bones, the sign of cunning and cupidity, or his pointed chin, the sign of acuteness and perseverance. Planchet reigned with as much majesty in his dining-room as in his shop. He set before his master a frugal, but perfectly Parisian repast: roast meat, cooked at the baker's, with vegetables, salad, and a dessert borrowed from the shop itself. D'Artagnan was pleased that the grocer had drawn from behind the fagots a bottle of that Anjou wine which during all his life had been D'Artagnan's favourite wine.

<Formerly, monsieur,> said Planchet, with a smile full of bonhomie, <it was I who drank your wine; now you do me the honour to drink mine.>

<And, thank God, friend Planchet, I shall drink it for a long time to come, I hope; for at present I am free.>

<Free? You have a leave of absence, monsieur?>

<Unlimited.>

<You are leaving the service?> said Planchet, stupefied.

<Yes, I am resting.>

<And the king?> cried Planchet, who could not suppose it possible that the king could do without the services of such a man as D'Artagnan.

<The king will try his fortune elsewhere. But we have supped well, you are disposed to enjoy yourself; you invite me to confide in you. Open your ears, then.>

<They are open.> And Planchet, with a laugh more frank than cunning, opened a bottle of white wine.

<Leave me my reason, at least.>

<Oh, as to you losing your head—you, monsieur!>

<Now my head is my own, and I mean to take better care of it than ever. In the first place we shall talk business. How fares our money-box?>

<Wonderfully well, monsieur. The twenty thousand livres I had of you are still employed in my trade, in which they bring me nine per cent. I give you seven, so I gain two by you.>

<And you are still satisfied?>

<Delighted. Have you brought me any more?>

<Better than that. But do you want any?>

<Oh! not at all. Every one is willing to trust me now. I am extending my business.>

<That was your intention.>

<I play the banker a little. I buy goods of my needy brethren; I lend money to those who are not ready for their payments.>

<Without usury?>

<Oh! monsieur, in the course of the last week I have had two meetings on the boulevards, on account of the word you have just pronounced.>

<What?>

<You shall see: it concerned a loan. The borrower gives me in pledge some raw sugars, on condition that I should sell if repayment were not made within a fixed period. I lend a thousand livres. He does not pay me, and I sell the sugars for thirteen hundred livres. He learns this and claims a hundred crowns. \swear{Ma foi!} I refused, pretending that I could not sell them for more than nine hundred livres. He accused me of usury. I begged him to repeat that word to me behind the boulevards. He was an old guard, and he came: and I passed your sword through his left thigh.>

<\swear{Tu dieu!} what a pretty sort of banker you make!> said D'Artagnan.

<For above thirteen per cent I fight,> replied Planchet; <that is my character.>

<Take only twelve,> said D'Artagnan, <and call the rest premium and brokerage.>

<You are right, monsieur; but to your business.>

<Ah! Planchet, it is very long and very hard to speak.>

<Do speak it, nevertheless.>

D'Artagnan twisted his moustache like a man embarrassed with the confidence he is about to make and mistrustful of his confidant.

<Is it an investment?> asked Planchet.

<Why, yes.>

<At good profit?>

<A capital profit,—four hundred per cent, Planchet.>

Planchet gave such a blow with his fist upon the table, that the bottles bounded as if they had been frightened.

<Good heavens! is that possible?>

<I think it will be more,> replied D'Artagnan coolly; <but I like to lay it at the lowest!>

<The devil!> said Planchet, drawing nearer. <Why, monsieur, that is magnificent! Can one put much money in it?>

<Twenty thousand livres each, Planchet.>

<Why, that is all you have, monsieur. For how long a time?>

<For a month.>

<And that will give us\longdash>

<Fifty thousand livres each, profit.>

<It is monstrous! It is worth while to fight for such interest as that!>

<In fact, I believe it will be necessary to fight not a little,> said D'Artagnan, with the same tranquillity; <but this time there are two of us, Planchet, and I shall take all the blows to myself.>

<Oh! monsieur, I will not allow that.>

<Planchet, you cannot be concerned in it; you would be obliged to leave your business and your family.>

<The affair is not in Paris, then.>

<No.>

<Abroad?>

<In England.>

<A speculative country, that is true,> said Planchet,—<a country that I know well. What sort of an affair, monsieur, without too much curiosity?>

<Planchet, it is a restoration.>

<Of monuments?>

<Yes, of monuments; we shall restore Whitehall.>

<That is important. And in a month, you think?>

<I shall undertake it.>

<That concerns you, monsieur, and when once you are engaged\longdash>

<Yes, that concerns me. I know what I am about; nevertheless, I will freely consult with you.>

<You do me great honour; but I know very little about architecture.>

<Planchet, you are wrong; you are an excellent architect, quite as good as I am, for the case in question.>

<Thanks, monsieur. But your old friends of the musketeers?>

<I have been, I confess, tempted to speak of the thing to those gentlemen, but they are all absent from their houses. It is vexatious, for I know none more bold or able.>

<Ah! then it appears there will be an opposition, and the enterprise will be disputed?>

<Oh, yes, Planchet, yes.>

<I burn to know the details, monsieur.>

<Here they are, Planchet—close all the doors tight.>

<Yes, monsieur.> And Planchet double-locked them.

<That is well; now draw near.> Planchet obeyed.

<And open the window, because the noise of the passers-by and the carts will deafen all who might hear us.> Planchet opened the window as desired, and the gust of tumult which filled the chamber with cries, wheels, barkings, and steps deafened D'Artagnan himself, as he had wished. He then swallowed a glass of white wine, and began in these terms: <Planchet, I have an idea.>

<Ah! monsieur, I recognize you so well in that!> replied Planchet, panting with emotion. 